% --- Archon physics formalization source begin ---
% archon:physics
% archon:covers PhyXMiniProblems/problem_phyx_mini_0470.lean
% archon:source-report reports/phyx_mini/problem_phyx_mini_0470.source.json

\chapter{Physics problem phyx\_mini\_0470}
\label{ch:PhyXMiniProblems_problem_phyx_mini_0470}

\paragraph{Problem source.}
\#\# Physical scenario

Figure shows a pV-diagram for a cyclic process in which the initial and final states of some thermodynamic system are the same. The state of the system starts at point a and proceeds counterclockwise in the pV-diagram to point b, then back to a; the total work is W = -500 J.

\#\# Question

Find the heat added during this process.

\#\# Answer choices

- A: A: 350J

- B: B: 500J

- C: C: 333J

- D: D: 450J

\#\# Figure caption (auxiliary; use the image as primary evidence)

The figure is a pressure-volume (\textbackslash{}(p\textbackslash{})-\textbackslash{}(V\textbackslash{})) diagram. Key elements include:

- **Axes**: The horizontal axis represents volume (\textbackslash{}(V\textbackslash{})) and the vertical axis represents pressure (\textbackslash{}(p\textbackslash{})).
- **Points and Lines**: 
  - Two points labeled \textbackslash{}(a\textbackslash{}) and \textbackslash{}(b\textbackslash{}) are marked on the curve. 
  - \textbackslash{}(a\textbackslash{}) is at pressure \textbackslash{}(p\_a\textbackslash{}) and volume \textbackslash{}(V\_a\textbackslash{}).
  - \textbackslash{}(b\textbackslash{}) is at pressure \textbackslash{}(p\_b\textbackslash{}) and volume \textbackslash{}(V\_b\textbackslash{}).
  - A curved line connects points \textbackslash{}(a\textbackslash{}) and \textbackslash{}(b\textbackslash{}), indicating a process path.
- **Shaded Area**: 
  - The area under the process curve between \textbackslash{}(a\textbackslash{}) and \textbackslash{}(b\textbackslash{}) is shaded and striped, indicating the work done by the system in this cycle. This area is enclosed by the path.
  - The area in the diagram suggests a cyclic process \textbackslash{}(a \textbackslash{}to b \textbackslash{}to a\textbackslash{}).
- **Text Annotation**: 
  - The annotation states, "Area enclosed by path = total work \textbackslash{}(W\textbackslash{}) done by system in process \textbackslash{}(a \textbackslash{}to b \textbackslash{}to a\textbackslash{}). In this case \textbackslash{}(W < 0\textbackslash{})."
  
The arrows suggest the direction of the process, indicating that the process goes from \textbackslash{}(a\textbackslash{}) to \textbackslash{}(b\textbackslash{}) and back to \textbackslash{}(a\textbackslash{}).

\#\# Dataset metadata

Laws of Thermodynamics; Physical Model Grounding Reasoning, Multi-Formula Reasoning

\paragraph{Recorded answer/context.}
B: B: 500J

\paragraph{Figure/image path.}
/root/proposal\_for\_physic/science-mango/phyx\_mini\_run/phyx\_data/test\_image/470.png

\paragraph{Formalization target.}
create a compiling Lean file with sorry bodies at `PhyXMiniProblems/problem\_phyx\_mini\_0470.lean`. The Lean declarations must preserve the physical quantities, dimensions or dimensional roles, figure labels, governing-law hypotheses, and final relation expressed by this problem.
Use Mathlib/Physlib names found through LeanExplore where available. If a physics API is missing, introduce faithful local abstractions rather than scalar placeholder aliases.

\begin{theorem}[Physics formalization target]
\label{thm:physics:phyx_mini_0470:target}
\lean{PhyXMiniProblems.ProblemPhyXMini0470.heatTransferDuringCycle_matches_recordedChoice}
\uses{def:physics:phyx-mini-0470:phyxminiproblems-problemphyxmini0470-energyinjoules, def:physics:phyx-mini-0470:phyxminiproblems-problemphyxmini0470-cyclicpvprocess, def:physics:phyx-mini-0470:phyxminiproblems-problemphyxmini0470-matchesproblemstatement, def:physics:phyx-mini-0470:phyxminiproblems-problemphyxmini0470-matchesprimarypressurevolumediagram, def:physics:phyx-mini-0470:phyxminiproblems-problemphyxmini0470-hasphysicalstatereadouts, def:physics:phyx-mini-0470:phyxminiproblems-problemphyxmini0470-satisfiespressurevolumeworkarealaw, def:physics:phyx-mini-0470:phyxminiproblems-problemphyxmini0470-satisfiesclosedsystemfirstlaw, def:physics:phyx-mini-0470:phyxminiproblems-problemphyxmini0470-netheattransfermagnitudeinjoules, def:physics:phyx-mini-0470:phyxminiproblems-problemphyxmini0470-displayedheatamountinjoules, def:physics:phyx-mini-0470:phyxminiproblems-problemphyxmini0470-recordedanswerchoice}
The assigned autoformalize agent should translate this physics problem into Lean declarations in the covered file, with theorem and lemma proof bodies written as `by sorry`.
\end{theorem}
\begin{proof}
This is an autoformalization task, not a proof task. Produce statements that can later be proved by the physics prover without weakening the physical model.
\end{proof}

% --- Archon declaration topology begin ---
\section{Declaration topology}
\begin{definition}[Volume Quantity]
  \label{def:physics:phyx-mini-0470:phyxminiproblems-problemphyxmini0470-volumequantity}
  \lean{PhyXMiniProblems.ProblemPhyXMini0470.VolumeQuantity}
  A nonnegative physical volume, carrying dimension L³
\end{definition}
\begin{proof}
  This is the declared model definition of Volume Quantity.
\end{proof}
\begin{definition}[Pressure Quantity]
  \label{def:physics:phyx-mini-0470:phyxminiproblems-problemphyxmini0470-pressurequantity}
  \lean{PhyXMiniProblems.ProblemPhyXMini0470.PressureQuantity}
  A physical pressure, represented by Physlib's dimensionful pressure
\end{definition}
\begin{proof}
  This is the declared model definition of Pressure Quantity.
\end{proof}
\begin{definition}[Energy Quantity]
  \label{def:physics:phyx-mini-0470:phyxminiproblems-problemphyxmini0470-energyquantity}
  \lean{PhyXMiniProblems.ProblemPhyXMini0470.EnergyQuantity}
  A signed physical energy, used for internal energy, heat, and work
\end{definition}
\begin{proof}
  This is the declared model definition of Energy Quantity.
\end{proof}
\begin{definition}[pressure In Pascals]
  \label{def:physics:phyx-mini-0470:phyxminiproblems-problemphyxmini0470-pressureinpascals}
  \lean{PhyXMiniProblems.ProblemPhyXMini0470.pressureInPascals}
  \uses{def:physics:phyx-mini-0470:phyxminiproblems-problemphyxmini0470-pressurequantity}
  Read a physical pressure in coherent SI pascals
\end{definition}
\begin{proof}
  This is the declared model definition of pressure In Pascals.
\end{proof}
\begin{definition}[volume In Cubic Meters]
  \label{def:physics:phyx-mini-0470:phyxminiproblems-problemphyxmini0470-volumeincubicmeters}
  \lean{PhyXMiniProblems.ProblemPhyXMini0470.volumeInCubicMeters}
  \uses{def:physics:phyx-mini-0470:phyxminiproblems-problemphyxmini0470-volumequantity}
  Read a physical volume in coherent SI cubic metres
\end{definition}
\begin{proof}
  This is the declared model definition of volume In Cubic Meters.
\end{proof}
\begin{definition}[energy In Joules]
  \label{def:physics:phyx-mini-0470:phyxminiproblems-problemphyxmini0470-energyinjoules}
  \lean{PhyXMiniProblems.ProblemPhyXMini0470.energyInJoules}
  \uses{def:physics:phyx-mini-0470:phyxminiproblems-problemphyxmini0470-energyquantity}
  Read signed dimensionful energy in coherent SI joules
\end{definition}
\begin{proof}
  This is the declared model definition of energy In Joules.
\end{proof}
\begin{definition}[Diagram Point]
  \label{def:physics:phyx-mini-0470:phyxminiproblems-problemphyxmini0470-diagrampoint}
  \lean{PhyXMiniProblems.ProblemPhyXMini0470.DiagramPoint}
  The two labelled equilibrium points in the supplied diagram
\end{definition}
\begin{proof}
  This is the declared model definition of Diagram Point.
\end{proof}
\begin{definition}[Process Stage]
  \label{def:physics:phyx-mini-0470:phyxminiproblems-problemphyxmini0470-processstage}
  \lean{PhyXMiniProblems.ProblemPhyXMini0470.ProcessStage}
  Stages used to distinguish the initial and final visits to state a
\end{definition}
\begin{proof}
  This is the declared model definition of Process Stage.
\end{proof}
\begin{definition}[Cycle Leg]
  \label{def:physics:phyx-mini-0470:phyxminiproblems-problemphyxmini0470-cycleleg}
  \lean{PhyXMiniProblems.ProblemPhyXMini0470.CycleLeg}
  The two directed legs of the closed process
\end{definition}
\begin{proof}
  This is the declared model definition of Cycle Leg.
\end{proof}
\begin{definition}[Axis Quantity]
  \label{def:physics:phyx-mini-0470:phyxminiproblems-problemphyxmini0470-axisquantity}
  \lean{PhyXMiniProblems.ProblemPhyXMini0470.AxisQuantity}
  Physical quantity assigned to one axis of the plot
\end{definition}
\begin{proof}
  This is the declared model definition of Axis Quantity.
\end{proof}
\begin{definition}[Curve Location]
  \label{def:physics:phyx-mini-0470:phyxminiproblems-problemphyxmini0470-curvelocation}
  \lean{PhyXMiniProblems.ProblemPhyXMini0470.CurveLocation}
  Qualitative location of one curved leg in the primary image
\end{definition}
\begin{proof}
  This is the declared model definition of Curve Location.
\end{proof}
\begin{definition}[Cycle Orientation]
  \label{def:physics:phyx-mini-0470:phyxminiproblems-problemphyxmini0470-cycleorientation}
  \lean{PhyXMiniProblems.ProblemPhyXMini0470.CycleOrientation}
  Orientation of a closed path in the pressure--volume plane
\end{definition}
\begin{proof}
  This is the declared model definition of Cycle Orientation.
\end{proof}
\begin{definition}[Work Sign]
  \label{def:physics:phyx-mini-0470:phyxminiproblems-problemphyxmini0470-worksign}
  \lean{PhyXMiniProblems.ProblemPhyXMini0470.WorkSign}
  Sign displayed for the work done by the system
\end{definition}
\begin{proof}
  This is the declared model definition of Work Sign.
\end{proof}
\begin{definition}[Enclosed Area Meaning]
  \label{def:physics:phyx-mini-0470:phyxminiproblems-problemphyxmini0470-enclosedareameaning}
  \lean{PhyXMiniProblems.ProblemPhyXMini0470.EnclosedAreaMeaning}
  Physical meaning assigned to the shaded enclosed area by the annotation
\end{definition}
\begin{proof}
  This is the declared model definition of Enclosed Area Meaning.
\end{proof}
\begin{definition}[Thermodynamic State]
  \label{def:physics:phyx-mini-0470:phyxminiproblems-problemphyxmini0470-thermodynamicstate}
  \lean{PhyXMiniProblems.ProblemPhyXMini0470.ThermodynamicState}
  \uses{def:physics:phyx-mini-0470:phyxminiproblems-problemphyxmini0470-volumequantity, def:physics:phyx-mini-0470:phyxminiproblems-problemphyxmini0470-pressurequantity, def:physics:phyx-mini-0470:phyxminiproblems-problemphyxmini0470-energyquantity}
  Pressure, volume, and internal energy at one equilibrium stage
\end{definition}
\begin{proof}
  This is the declared model definition of Thermodynamic State.
\end{proof}
\begin{definition}[Pressure Volume Diagram]
  \label{def:physics:phyx-mini-0470:phyxminiproblems-problemphyxmini0470-pressurevolumediagram}
  \lean{PhyXMiniProblems.ProblemPhyXMini0470.PressureVolumeDiagram}
  \uses{def:physics:phyx-mini-0470:phyxminiproblems-problemphyxmini0470-volumequantity, def:physics:phyx-mini-0470:phyxminiproblems-problemphyxmini0470-pressurequantity, def:physics:phyx-mini-0470:phyxminiproblems-problemphyxmini0470-diagrampoint, def:physics:phyx-mini-0470:phyxminiproblems-problemphyxmini0470-cycleleg, def:physics:phyx-mini-0470:phyxminiproblems-problemphyxmini0470-axisquantity, def:physics:phyx-mini-0470:phyxminiproblems-problemphyxmini0470-curvelocation, def:physics:phyx-mini-0470:phyxminiproblems-problemphyxmini0470-cycleorientation, def:physics:phyx-mini-0470:phyxminiproblems-problemphyxmini0470-worksign, def:physics:phyx-mini-0470:phyxminiproblems-problemphyxmini0470-enclosedareameaning}
  Literal and calibrated data belonging to the supplied p--V image. The signed enclosed-area readout has energy units because pressure times volume has the dimension of work
\end{definition}
\begin{proof}
  This is the declared model definition of Pressure Volume Diagram.
\end{proof}
\begin{definition}[Cyclic PVProcess]
  \label{def:physics:phyx-mini-0470:phyxminiproblems-problemphyxmini0470-cyclicpvprocess}
  \lean{PhyXMiniProblems.ProblemPhyXMini0470.CyclicPVProcess}
  \uses{def:physics:phyx-mini-0470:phyxminiproblems-problemphyxmini0470-energyquantity, def:physics:phyx-mini-0470:phyxminiproblems-problemphyxmini0470-diagrampoint, def:physics:phyx-mini-0470:phyxminiproblems-problemphyxmini0470-processstage, def:physics:phyx-mini-0470:phyxminiproblems-problemphyxmini0470-thermodynamicstate, def:physics:phyx-mini-0470:phyxminiproblems-problemphyxmini0470-pressurevolumediagram}
  Independent observables for the complete cycle. In particular, net heat is an unconstrained dimensionful field here; it is related to work and internal energy only by the governing first-law hypothesis below
\end{definition}
\begin{proof}
  This is the declared model definition of Cyclic PVProcess.
\end{proof}
\begin{definition}[Matches Problem Statement]
  \label{def:physics:phyx-mini-0470:phyxminiproblems-problemphyxmini0470-matchesproblemstatement}
  \lean{PhyXMiniProblems.ProblemPhyXMini0470.MatchesProblemStatement}
  \uses{def:physics:phyx-mini-0470:phyxminiproblems-problemphyxmini0470-energyinjoules, def:physics:phyx-mini-0470:phyxminiproblems-problemphyxmini0470-cyclicpvprocess}
  Facts supplied by the prose: the route starts at a, visits b, returns to the same thermodynamic state a, and has total work -500 J. No heat value is assumed
\end{definition}
\begin{proof}
  This is the declared model definition of Matches Problem Statement.
\end{proof}
\begin{definition}[Matches Primary Pressure Volume Diagram]
  \label{def:physics:phyx-mini-0470:phyxminiproblems-problemphyxmini0470-matchesprimarypressurevolumediagram}
  \lean{PhyXMiniProblems.ProblemPhyXMini0470.MatchesPrimaryPressureVolumeDiagram}
  \uses{def:physics:phyx-mini-0470:phyxminiproblems-problemphyxmini0470-pressureinpascals, def:physics:phyx-mini-0470:phyxminiproblems-problemphyxmini0470-volumeincubicmeters, def:physics:phyx-mini-0470:phyxminiproblems-problemphyxmini0470-diagrampoint, def:physics:phyx-mini-0470:phyxminiproblems-problemphyxmini0470-cyclicpvprocess}
  Primary-image evidence: volume is horizontal, pressure is vertical, V\_a < V\_b, p\_b < p\_a, both states and labels are shown, and the arrows trace the lower a → b curve followed by the upper b → a curve. The annotation identifies the signed enclosed area with negative work by the system, but does not assign any heat
\end{definition}
\begin{proof}
  This is the declared model definition of Matches Primary Pressure Volume Diagram.
\end{proof}
\begin{definition}[Has Physical State Readouts]
  \label{def:physics:phyx-mini-0470:phyxminiproblems-problemphyxmini0470-hasphysicalstatereadouts}
  \lean{PhyXMiniProblems.ProblemPhyXMini0470.HasPhysicalStateReadouts}
  \uses{def:physics:phyx-mini-0470:phyxminiproblems-problemphyxmini0470-pressureinpascals, def:physics:phyx-mini-0470:phyxminiproblems-problemphyxmini0470-volumeincubicmeters, def:physics:phyx-mini-0470:phyxminiproblems-problemphyxmini0470-cyclicpvprocess}
  Positivity conditions selecting genuine pressure--volume states
\end{definition}
\begin{proof}
  This is the declared model definition of Has Physical State Readouts.
\end{proof}
\begin{definition}[Satisfies Pressure Volume Work Area Law]
  \label{def:physics:phyx-mini-0470:phyxminiproblems-problemphyxmini0470-satisfiespressurevolumeworkarealaw}
  \lean{PhyXMiniProblems.ProblemPhyXMini0470.SatisfiesPressureVolumeWorkAreaLaw}
  \uses{def:physics:phyx-mini-0470:phyxminiproblems-problemphyxmini0470-energyinjoules, def:physics:phyx-mini-0470:phyxminiproblems-problemphyxmini0470-cyclicpvprocess}
  The quasi-static p--V work law for the complete closed path: its signed enclosed area equals the net work done by the system. This is a governing relation, not the requested heat conclusion
\end{definition}
\begin{proof}
  This is the declared model definition of Satisfies Pressure Volume Work Area Law.
\end{proof}
\begin{definition}[Satisfies Closed System First Law]
  \label{def:physics:phyx-mini-0470:phyxminiproblems-problemphyxmini0470-satisfiesclosedsystemfirstlaw}
  \lean{PhyXMiniProblems.ProblemPhyXMini0470.SatisfiesClosedSystemFirstLaw}
  \uses{def:physics:phyx-mini-0470:phyxminiproblems-problemphyxmini0470-energyinjoules, def:physics:phyx-mini-0470:phyxminiproblems-problemphyxmini0470-cyclicpvprocess}
  Closed-system thermodynamics for the complete traversal. Internal energy is a state function, and with heat positive into the system and work positive by the system the first law is Q\_into = ΔU + W\_by
\end{definition}
\begin{proof}
  This is the declared model definition of Satisfies Closed System First Law.
\end{proof}
\begin{definition}[net Heat Transfer Magnitude In Joules]
  \label{def:physics:phyx-mini-0470:phyxminiproblems-problemphyxmini0470-netheattransfermagnitudeinjoules}
  \lean{PhyXMiniProblems.ProblemPhyXMini0470.netHeatTransferMagnitudeInJoules}
  \uses{def:physics:phyx-mini-0470:phyxminiproblems-problemphyxmini0470-energyinjoules, def:physics:phyx-mini-0470:phyxminiproblems-problemphyxmini0470-cyclicpvprocess}
  Magnitude of the net heat transfer, in joules
\end{definition}
\begin{proof}
  This is the declared model definition of net Heat Transfer Magnitude In Joules.
\end{proof}
\begin{definition}[Answer Choice]
  \label{def:physics:phyx-mini-0470:phyxminiproblems-problemphyxmini0470-answerchoice}
  \lean{PhyXMiniProblems.ProblemPhyXMini0470.AnswerChoice}
  Labels of the four displayed multiple-choice answers
\end{definition}
\begin{proof}
  This is the declared model definition of Answer Choice.
\end{proof}
\begin{definition}[displayed Heat Amount In Joules]
  \label{def:physics:phyx-mini-0470:phyxminiproblems-problemphyxmini0470-displayedheatamountinjoules}
  \lean{PhyXMiniProblems.ProblemPhyXMini0470.displayedHeatAmountInJoules}
  \uses{def:physics:phyx-mini-0470:phyxminiproblems-problemphyxmini0470-answerchoice}
  Positive energy amount printed beside each choice, in joules
\end{definition}
\begin{proof}
  This is the declared model definition of displayed Heat Amount In Joules.
\end{proof}
\begin{definition}[recorded Answer Choice]
  \label{def:physics:phyx-mini-0470:phyxminiproblems-problemphyxmini0470-recordedanswerchoice}
  \lean{PhyXMiniProblems.ProblemPhyXMini0470.recordedAnswerChoice}
  \uses{def:physics:phyx-mini-0470:phyxminiproblems-problemphyxmini0470-answerchoice}
  The dataset records answer label B
\end{definition}
\begin{proof}
  This is the declared model definition of recorded Answer Choice.
\end{proof}
\begin{lemma}[net Internal Energy Change In Joules eq zero]
  \label{lem:physics:phyx-mini-0470:phyxminiproblems-problemphyxmini0470-netinternalenergychangeinjoules-eq-zero}
  \lean{PhyXMiniProblems.ProblemPhyXMini0470.netInternalEnergyChangeInJoules_eq_zero}
  \uses{def:physics:phyx-mini-0470:phyxminiproblems-problemphyxmini0470-energyinjoules, def:physics:phyx-mini-0470:phyxminiproblems-problemphyxmini0470-cyclicpvprocess, def:physics:phyx-mini-0470:phyxminiproblems-problemphyxmini0470-matchesproblemstatement, def:physics:phyx-mini-0470:phyxminiproblems-problemphyxmini0470-satisfiesclosedsystemfirstlaw}
  Because the initial and final thermodynamic states coincide, their internal energies coincide and the net internal-energy change is zero
\end{lemma}
\begin{proof}
  The conclusion follows from the governing relations and figure readouts named in the statement of net Internal Energy Change In Joules eq zero.
\end{proof}
% --- Archon declaration topology end ---
% --- Archon physics formalization source end ---
